\section{The Angular Momentum Operator}

In classical mechanics, the angular momentum of a particle with position vector
$\mathbf{r}$ and linear momentum $\mathbf{p}$ is defined as the vector product
$\mathbf{L}_{cl} = \mathbf{r} \times \mathbf{p}$. To transition to quantum
mechanics, we apply the correspondence principle, promoting the dynamical
variables to linear operators acting on the Hilbert space: $\mathbf{r} \to
\hat{\mathbf{r}}$ and $\mathbf{p} \to \hat{\mathbf{p}}$.

However, since $\hat{\mathbf{r}}$ and $\hat{\mathbf{p}}$ are non-commuting
operators, the definition of the vector product must be handled with algebraic
rigor to ensure the resulting operator is Hermitian (measurable) and
unambiguous.

\subsection{Standard Vector Definition}

In 3D Euclidean space, we define the operator $\hat{\mathbf{L}}$ as the vector
product of the position and momentum operators:
\begin{equation}
  \hat{\mathbf{L}} \equiv \hat{\mathbf{r}} \times \hat{\mathbf{p}} = -i\hbar
  (\mathbf{r} \times \nabla).
\end{equation}
Using the Levi-Civita pseudotensor $\epsilon_{ijk}$, the components are defined
as:
\begin{equation} \label{eq:L_tensor_def}
  \hat{L}_i = \epsilon_{ijk} \hat{x}_j \hat{p}_k,
\end{equation}
where Einstein summation is implied. This definition ensures that
$\hat{\mathbf{L}}$ acts as a vector operator under spatial rotations.

\subsection{Geometric Algebra Perspective (Clifford Algebra)}

To provide a coordinate-independent definition valid in any dimension, we employ
the Geometric Product within the algebra $\mathcal{C}\ell_{3,0}(\mathbb{R})$.
The product of the vector operators $\hat{\mathbf{r}}$ and $\hat{\mathbf{p}}$
decomposes into a scalar (symmetric) and bivector (antisymmetric) part:
\begin{equation}
  \hat{\mathbf{r}}\hat{\mathbf{p}} = \hat{\mathbf{r}} \cdot \hat{\mathbf{p}} +
  \hat{\mathbf{r}} \wedge \hat{\mathbf{p}}.
\end{equation}
The physically significant quantity is the bivector $\hat{\mathcal{L}}$,
representing the oriented plane of rotation:
\begin{equation}
  \hat{\mathcal{L}} \equiv \hat{\mathbf{r}} \wedge \hat{\mathbf{p}} =
  \frac{1}{2}(\hat{\mathbf{r}}\hat{\mathbf{p}} -
  \hat{\mathbf{p}}\hat{\mathbf{r}}).
\end{equation}
The standard vector operator $\hat{\mathbf{L}}$ is recovered as the Hodge
dual of this bivector. Using the unit pseudoscalar $I = \hat{e}_1 \hat{e}_2
\hat{e}_3$:
\begin{equation}
  \hat{\mathbf{L}} = -i I \hat{\mathcal{L}}.
\end{equation}
This formulation reveals that spin and orbital angular momentum share the same
bivector structure, differing only in the vector space they act upon.

\subsection{Hermiticity and Observability}

For $\hat{\mathbf{L}}$ to represent a physical observable, it must be Hermitian.
We examine the adjoint of the component definition (\ref{eq:L_tensor_def}):
\begin{equation}
  \hat{L}_i^\dagger = (\epsilon_{ijk} \hat{x}_j \hat{p}_k)^\dagger.
\end{equation}
Using the property $(\hat{A}\hat{B})^\dagger = \hat{B}^\dagger \hat{A}^\dagger$
and the fact that $\hat{\mathbf{x}}$ and $\hat{\mathbf{p}}$ are Hermitian
($\hat{x}_j^\dagger = \hat{x}_j$, $\hat{p}_k^\dagger = \hat{p}_k$) and
$\epsilon_{ijk}$ is real:
\begin{equation}
  \hat{L}_i^\dagger = \epsilon_{ijk} \hat{p}_k^\dagger \hat{x}_j^\dagger =
  \epsilon_{ijk} \hat{p}_k \hat{x}_j.
\end{equation}
The non-commutativity of $\hat{\mathbf{x}}$ and $\hat{\mathbf{p}}$ is resolved
by the canonical commutation relation $[\hat{x}_j, \hat{p}_k] = i\hbar
\delta_{jk}$:
\begin{equation}
  [\hat{x}_j, \hat{p}_k] = i\hbar \delta_{jk} \implies \hat{p}_k \hat{x}_j =
  \hat{x}_j \hat{p}_k - i\hbar \delta_{jk}.
\end{equation}
Substituting this back into the expression for $\hat{L}_i^\dagger$:
\begin{align}
  \hat{L}_i^\dagger &= \epsilon_{ijk} (\hat{x}_j \hat{p}_k - i\hbar \delta_{jk})
  \nonumber \\
                    &= \epsilon_{ijk} \hat{x}_j \hat{p}_k - i\hbar \epsilon_{ijk}
  \delta_{jk}.
\end{align}
The second term contains the contraction of the antisymmetric tensor
$\epsilon_{ijk}$ with the symmetric Kronecker delta $\delta_{jk}$. Since
$\epsilon_{ijk}$ vanishes if any two indices are equal (which $\delta_{jk}$
forces), we have:
\begin{equation}
  \epsilon_{ijk} \delta_{jk} = \epsilon_{ijj} = 0.
\end{equation}
Therefore:
\begin{equation}
  \hat{L}_i^\dagger = \epsilon_{ijk} \hat{x}_j \hat{p}_k = \hat{L}_i.
\end{equation}
The operator is inherently Hermitian, and no symmetrization (e.g.,
$\frac{1}{2}(\mathbf{r}\times\mathbf{p} - \mathbf{p}\times\mathbf{r})$) is
required because the cross product couples only commuting orthogonal components
(e.g., $x$ and $p_y$).

\subsection{Explicit Component Form}

Expanding the tensor notation, we obtain the explicit forms used in standard
calculations:

\begin{align}
  \hat{L}_x &= \hat{y}\hat{p}_z - \hat{z}\hat{p}_y, \\
  \hat{L}_y &= \hat{z}\hat{p}_x - \hat{x}\hat{p}_z, \\
  \hat{L}_z &= \hat{x}\hat{p}_y - \hat{y}\hat{p}_x.
\end{align}

These expressions confirm that the quantum angular momentum operator preserves
the functional form of its classical counterpart, provided the order of
operators is maintained as $\mathbf{r} \times \mathbf{p}$.

\subsection{The Lie Algebra \texorpdfstring{$\mathfrak{so}(3)$}{so(3)}}

The fundamental commutation relations for the components of $\hat{\mathbf{L}}$
define the structure of the rotation group. Computing $[\hat{L}_x, \hat{L}_y]$:
\begin{align}
  [\hat{L}_x, \hat{L}_y] &= [y p_z - z p_y, z p_x - x p_z] \nonumber \\
                         &= [y p_z, z p_x] + [-z p_y, -x p_z] \quad
                         (\text{other terms commute}) \nonumber \\
                         &= y p_x [p_z, z] + x p_y [z, p_z] \nonumber \\
                         &= y p_x (-i\hbar) + x p_y (i\hbar) \nonumber \\
                         &= i\hbar (x p_y - y p_x) = i\hbar \hat{L}_z.
\end{align}
Generalizing this result yields the Lie algebra of $\mathfrak{so}(3)$:
\begin{equation}
  [\hat{L}_i, \hat{L}_j] = i\hbar \epsilon_{ijk} \hat{L}_k.
\end{equation}
This closure relation confirms that $\hat{\mathbf{L}}$ generates the group of
spatial rotations $\mathrm{SO}(3)$.
